\begin{description}
\item[(a)] First its bolometric magnitude is calculated from equation:
  $M_V - m_V = 5 - 5\,\log(r)$ or equivalent:
  $M_V - m_V = 5 + 5\,\log\pi \;\rightarrow\;
  M_V = 12.2 + 5 + 5\log(0''.001) = 12.2 + 5 - 15 = 2.2$~mag.
  Its barycentric correction is: % sic: source says "barycentric" for bolometric
  $B.C. = M_{\mathrm{bol}} - M_V$ and $M_{\mathrm{bol}} = B.C. + M_V$ or
  $M_{\mathrm{bol}} = -0.6 + 2.2$ or $M_{\mathrm{bol}} = 1.6$~mag.
  \hfill(4 Points)

  Then its Luminosity is calculated from:
  \[
    M_\odot - M_{\mathrm{bol}} = 2.5\,\log\!\left(\frac{L}{L_\odot}\right),
    \ \text{or}\
    4.72 - 1.6 = 2.5\,\log\!\left(\frac{L}{L_\odot}\right)
    \ \text{or}\
    \log\!\left(\frac{L}{L_\odot}\right) = 1.25
    \ \text{and}\ L = 17.70\,L_\odot
  \]
  \hfill(2 Point)
\item[(b)] Type of star: A star with $M_{\mathrm{bol}} = 1.6$~mag,
  $L = 17.7\,L_\odot$ and $T_{\mathrm{eff}} = 4000$~K is much brighter and much
  cooler than the Sun (see \emph{Table of constants}). Therefore it is (i) a red
  giant star.
  \hfill(4 Points)
\end{description}
